\chapter{Introduction}
\label{ch:intro}

Spoken conversation is the most natural form of human--machine interaction.
Nevertheless, until recently spoken dialogue systems were almost always built
from a chain of independent components: voice activity detection, speech
recognition, a text dialogue engine, and then text-to-speech. The introduction of
models such as Moshi~\cite{defossez2024moshi} and
Freeze-Omni~\cite{wang2025freezeomni} showed that spoken dialogue can be modelled
as a single ``speech-to-speech'' problem with one integrated model: a model that
listens and speaks at the same time and needs no explicit speech segmentation for
turn-taking.

This shift has happened for high-resource languages such as English and Chinese.
For Persian, none of the open speech-to-speech models offers official support,
and Persian conversational datasets that cover overlapping speech, interruptions
and backchannels are not publicly available either. The practical question of
this project arises from this gap: with the computing and time budget of a
bachelor's project, how far can one go in building a full-duplex Persian
speech-to-speech system, and, more importantly, what can be ``proven'' about it
with auditable evidence?

\section{Problem Definition}
\label{sec:problem}

The defined goal of this project is to design, implement and evaluate a prototype
speech-to-speech large language model for Persian. The system must work in
\emph{full-duplex} mode: unlike half-duplex systems, which listen to the user only
after the reply has finished, the user's microphone must stay active while the
reply plays, and the system must be able to react if the user \emph{interrupts}.

Three properties separate this problem from a conventional spoken question
answering system:

\begin{itemize}
\item \textbf{Listening and speaking at the same time.} While playing its own
      audio output, the system must also process the microphone input stream and
      decide whether to continue or stop playback.
\item \textbf{Time-critical decisions.} The decision to stop playback must be made
      within a few tens of milliseconds, otherwise the user feels that the system
      does not hear them.
\item \textbf{Continuity of interrupted speech.} What the user says at the moment
      of the interruption must not be lost; it must be carried to the next turn,
      otherwise the user has to repeat the sentence.
\end{itemize}

\section{Motivation}
\label{sec:motivation}

Addressing this problem for Persian matters for several reasons.

First, \textbf{accessibility}. A speech interface can be the only effective way to
use digital services for users who are not comfortable with a Persian keyboard,
for users with low vision, and in particular for older adults. In the natural
conversation of older adults, longer pauses and mid-sentence corrections are more
frequent; a half-duplex system that does not listen until the end of its own
reply creates an irritating experience in these conditions.

Second, \textbf{Persian is low-resource in this particular area}. Persian speech
recognition and text-to-speech datasets are available; a prominent example is
ManaTTS~\cite{qharabagh2025manatts}, with more than 86 hours of single-speaker
speech. But these resources are single-speaker and read, and the interactive
phenomena that a full-duplex model has to learn --- overlap, interruptions and
backchannels --- do not occur in them.

Third, \textbf{computational cost}. Current speech-to-speech models, such as
Qwen2.5-Omni~\cite{xu2025qwen25omni}, need far more memory to run than a
laboratory graphics card. Showing what is ``achievable'' and what is ``not
achievable'' within a memory range of 12--24 GB is itself a useful engineering
result.

\section{Defined Requirements}
\label{sec:requirements}

The project definition sets seven specific requirements. They are used throughout
this report as the yardstick of achievement and are listed in
Table~\ref{tab:requirements}. For each requirement, the results chapter
(Chapter~\ref{ch:results}) reports the final status and the class of the
supporting evidence.

\begin{table}[ht]
\centering
\caption{The seven requirements of the project definition}
\label{tab:requirements}
\small
\begin{tabular}{clL{6.4cm}}
\toprule
\textbf{ID} & \textbf{Requirement} & \textbf{Specified measure} \\
\midrule
R1 & Persian speech-to-speech prototype & Persian spoken reply generated from
      Persian spoken input \\
R2 & Full-duplex operation & Listening during playback and reacting to
      interruptions \\
R3 & End-to-end latency & At most 500 ms from receiving speech to producing the
      spoken reply \\
R4 & Persian conversational dataset & 100--200 hours of speech with noise,
      overlap and interruptions \\
R5 & Interruption detection module & Accuracy above 80\% with energy, pitch and
      MFCC features \\
R6 & Evaluation on target hardware & A GPU with 12--24 GB of memory \\
R7 & Human evaluation & 5--10 Persian speakers, at least two aged over 60 \\
\bottomrule
\end{tabular}
\end{table}

The expected final outputs comprise four items: a working prototype, documented
source code, a Persian conversational dataset, and findings on how older
Persian-speaking adults interact with the system. For the fourth item, the
measurement protocol and infrastructure are ready; its empirical finding can only
be claimed once the human study is complete, so this report makes no claim about
the experience of older adults.

\section{Research Questions}
\label{sec:rq}

The project proposal poses four research questions, each of which
Chapter~\ref{ch:conclusions} (Section~\ref{sec:rq-answers}) answers with the
evidence of this report:

\begin{enumerate}
\item \textbf{Question 1.} How can an existing English or Chinese speech-to-speech
model be adapted to Persian with the least retraining cost?
\item \textbf{Question 2.} Which architectural changes, such as a dual buffer or
an interruption detection module, are needed for full-duplex operation with a
latency below 500 ms?
\item \textbf{Question 3.} What is the trade-off between latency, model size and
speech quality for Persian in real-time operation?
\item \textbf{Question 4.} How do older Persian-speaking adults perceive the
responsiveness and naturalness of such a system?
\end{enumerate}

\section{Main Challenges}
\label{sec:challenges}

Over the course of the project, four decisive challenges appeared that changed
the final shape of the work.

\paragraph{No official training path for Persian.}
Mini-Omni2~\cite{xie2024miniomni2}, which the project definition names as an
example base model, and LLaMA-Omni2~\cite{fang2025llamaomni2} have published
repositories for inference but provide no complete, official speech-to-speech
trainer. Among the models reviewed, only the Moshi family had published an
official low-rank adaptation trainer~\cite{moshifinetune2024}; for this reason
the direct track of this project was built on Moshika, the 7-billion-parameter
version of that family.

\paragraph{No labelled Persian conversational data.}
No public Persian corpus with interruption and overlap labels was available. The
dataset of this project therefore had to be built from an internal archive of
Persian spoken content with an automatic pipeline of window reconstruction,
speaker diarization and response-pair extraction. The consequence of this choice
is that the resulting interaction labels are ``automatic'' and not human-verified;
this boundary is kept throughout the report.

\paragraph{The gap between reducing loss and valid generation.}
This is the most important negative finding of the project. Over seven
experimental generations, low-rank adaptation of the direct model steadily
reduced the \emph{teacher-forced} loss, but never reached valid
\emph{autoregressive} generation of Persian speech. The final audit revealed its
structural cause: the model's English tokenizer needs more than six tokens per
Persian word and exceeds the budget of one token per frame.

\paragraph{The boundary between engineering evidence and scientific claims.}
Automatic measures were available, but 12--24 GB target hardware, human labelling
of interruptions and a large human study were not available within the project
deadline; instead, a browser loopback test with the memory budget of a 12 GB card,
independent two-pass automatic labels and a small listening study were designed.
The project's methodological decision was to label every measure with its
``evidence class'' and never to present an automatic approximation in place of
official evidence. Section~\ref{sec:evidence-first} turns this decision into a way
of working.

\section{Two-Track Strategy}
\label{sec:two-track}

To reduce execution risk, the project pursued two tracks in parallel:

\begin{enumerate}
\item \textbf{The direct track.} Low-rank adaptation of Moshika on Persian
      conversational data, in the hope of obtaining a single Persian
      speech-to-speech model. This track answers the spirit of the project
      definition directly.
\item \textbf{The cascade track.} A modular chain of Persian speech recognition,
      text reply generation with a large language model, and text-to-speech, on
      top of which full-duplex control and interruption logic are built
      separately. This track guarantees that the project has a working system to
      demonstrate.
\end{enumerate}

Both tracks used the same dataset. The final outcome, as reported in
Chapter~\ref{ch:results}, was asymmetric: the cascade track reached a streaming
full-duplex system that meets the numeric latency and interruption criteria in
automated tests, while the direct track, and a speech-input track with Gemma~4,
are reported as negative results together with their measured causes.

\section{Contributions}
\label{sec:contributions}

The contributions of this project can be summarized in six items.

\begin{enumerate}
\item \textbf{A streaming full-duplex Persian voice assistant.} The system works
      with an open microphone, detects the end of speech automatically and
      speculatively, and generates and plays the reply clause by clause. In
      loopback testing with a real browser, the mean latency from end of speech to
      the first audio of the reply was 372 ms and its 90th percentile 431 ms, no
      turn exceeded 500 ms, and the GPU side of the system fit in 11.1~GiB of
      memory.

\item \textbf{Two-stage interruption detection with lexical confirmation.} Ducking
      after one frame of speech louder than the expected echo, and confirmation
      with the speech recognizer after removing listener tokens and the
      assistant's own echo; accuracy 0.933 and F1 0.880 in a deployment
      simulation and accuracy 0.825 in a real browser, against 0.574 for an
      acoustic classifier alone; the reply after an interruption, keeping only the
      heard part of the interrupted reply in the history, was as relevant as the
      reply to a normal turn.

\item \textbf{A derived and audited Persian conversational dataset.} 6,754 ``user
      speech / response'' pairs with 108.584 hours of stereo audio, session-based
      group splits and zero group leakage; an independent automatic listener
      confirmed the turn boundary in 93.8\% of samples.

\item \textbf{Evidence-based component selection.} A final semantic test of the
      responder (relevance 3.150 versus 1.425), and re-selection of the recognizer
      on human transcripts, the responder on requests of older adults and the
      synthesizer with two automatic listeners, all under pre-registered
      protocols; and optimization of the lab synthesizer without retraining, which
      found and fixed several of its text front-end errors (a dropped final word,
      misread numbers).

\item \textbf{A measured cause for the failure of the direct track.} Seven
      experimental generations of Moshika and one speech-input model; it was shown
      that the tokenizer bottleneck is real and that removing it doubles the best
      result, but is not enough for intelligible speech.

\item \textbf{An evidence-first way of working.} Protocols and thresholds were
      registered before running, test splits were opened only once, result
      artifacts are write-once, and a \emph{fail-closed} aggregator prevents
      automatic approximations from being promoted to official evidence.
\end{enumerate}

The latency and interruption evidence is automated and on an H100, the labels are
not human, and the results of the listening study are not yet available. The exact
scope of these results is reported in Chapter~\ref{ch:discussion}.

\section{Structure of This Report}
\label{sec:structure}

This report is organized in seven chapters.

Chapter~\ref{ch:prelim} introduces the necessary preliminaries: the
representation of the speech signal and acoustic features, speech recognition and
text-to-speech, large language models and low-rank adaptation, neural audio
codecs and the architecture of speech-to-speech models, and the measures and
statistical concepts used in evaluation.

Chapter~\ref{ch:literature} reviews prior work and positions this project among
cascade systems and end-to-end spoken dialogue models, Persian speech resources,
and prior work on interruption detection.

Chapter~\ref{ch:method} describes the proposed method: the dataset construction
pipeline, the cascade architecture, the interruption detection module and the
full-duplex transport layer, the design of the direct-track experiments, and the
evidence-first way of working.

Chapter~\ref{ch:results} presents the results with the corresponding tables and
figures and states, for each result, its permitted claim boundary.

Chapter~\ref{ch:discussion} interprets the results, in particular the gap between
reducing loss and valid generation, and then lists the limitations and the
ethical and legal considerations.

Chapter~\ref{ch:conclusions} presents the conclusions and suggestions for future
work. The appendix gives the reproduction commands, the exact hyperparameters of
the experiments, and the map of evidence artifacts.
